%% —-----------------------------------------------------------
%% sbcreviews-pt.tex
%% Template para SBC Reviews 
%% (https://journals-sol.sbc.org.br/index.php/reviews/) 
%% 
%% Esta versão do template para a SBC Reviews inclui algumas
%% modificações sobre o template original da REIC
%% (https://journals-sol.sbc.org.br/index.php/reic)
%%
%% Licença CC BY 4.0 (https://creativecommons.org/licenses/by/4.0/)
%% Sociedade Brasileira de Computação.

\documentclass[portuguese]{sbcreviews-2025}%
\usepackage[misc,geometry]{ifsym} 

%\usepackage{graphicx} % carregada dentro da classe
%\usepackage{fontspec} % carregada dentro da classe
%\usepackage{fontawesome} % carregada dentro da classe
%\usepackage{color} % carregada dentro da classe
%\usepackage{hyperref} % carregada dentro da classe

\usepackage{aas_macros}
\usepackage[bottom]{footmisc}

\usepackage{tabularray}
% O pacote tabularray torna o uso dos pacotes abaixo desnecessarios.
%\usepackage{supertabular}
%\usepackage{multicol}
%\usepackage{multirow}

\usepackage{afterpage}
\usepackage{url}
\usepackage{pifont}

\setcitestyle{square}

\definecolor{engtitle}{rgb}{0.5,0.5,0.5}
\definecolor{orcidlogo}{rgb}{0.37,0.48,0.13}
\definecolor{unilogo}{rgb}{0.16, 0.26, 0.58}
\definecolor{maillogo}{rgb}{0.58, 0.16, 0.26}
\definecolor{darkblue}{rgb}{0.0,0.0,0.0}
\hypersetup{colorlinks,breaklinks,
            linkcolor=darkblue,urlcolor=darkblue,
            anchorcolor=darkblue,citecolor=darkblue}
%\hypersetup{colorlinks,citecolor=blue,linkcolor=blue,urlcolor=blue}
% We disable hyperlinks by default with this line, to avoid the error "\pdfendlink ended up in different nesting level" while writing.

\jid{ROCS}
\jtitle{SBC Computing Reviews, 202X, XX:1}
\issn{2966-3938}
\doi{10.5753/rocs.202X.XXXXXX}
\copyrightstatement{This work is licensed under a Creative Commons Attribution 4.0 International License}
\jyear{202X}


\category{Revisão da Literatura}
\title[Template da SBC Reviews 2025 para artigos em português]{Apresentação do template da SBC Reviews 2025 - Versão para artigos em português}
\engtitle{\textcolor{engtitle}{Presenting the SBC Reviews 2025 template - Version for papers in Portuguese}}

%THE ORCID IS MANDATORY FOR EACH AUTHOR
\author[Silva et al. 202X]{
\affil{\textbf{João Silva}~\orcidlink{0000-0000-0000-0000}~\textcolor{blue}{\faEnvelope}~~[{Universidade Federal}~|\href{mailto: ...}{~{\textit{...@...br}}}~]}

\affil{\textbf{Maria Santos}~\orcidlink{0000-0000-0000-0000}~[{Universidade Federal}~|\href{mailto: ...}{~{\textit{...@...br}}}~]}

\affil{\textbf{Ana Silva }~\orcidlink{0000-0000-0000-0000}~[{Universidade Federal}~|\href{mailto: ...}{~{\textit{...@...br}}}~]}

\affil{\textbf{José Santos}~\orcidlink{0000-0000-0000-0000}~[{Universidade Federal}~|\href{mailto: ...}{~{\textit{...@...br}}}~]}

}

\begin{document}

\begin{frontmatter}

\maketitle

\begin{mail}
Instituto ..., Universidade ..., Endereço, Cidade, Estado, CEP, Brasil. 
\end{mail}


\begin{abstract-pt}
Este texto, em formato de artigo científico, tem por objetivo apresentar o novo modelo para artigos da SBC, descrevendo suas principais características e explicando como deve ser utilizado. Esta versão, mais especificamente, deve ser utilizada exclusivamente para artigos escritos em português que serão publicados em alguma série de anais de eventos na SBC OpenLib. O resumo em português, como pode-se notar neste exemplo, deve vir antes do resumo em inglês (abstract) e ter entre 500 e 750 palavras.
\end{abstract-pt}

\begin{abstract-en}
This text, formatted as a scientific article, aims to present the new SBC paper template, describing its main features and explaining how it should be used. This version, more specifically, should be used exclusively for articles written in Portuguese that will be published in any event proceeding series at SBC OpenLib. The abstract in Portuguese, as you can see in this example, must be before the abstract in English and must have between 500 and 750 words.
\end{abstract-en}

\begin{pchaves}
Anais de evento, Modelo, SBC OpenLib, Indexação
\end{pchaves}

\begin{keywords}
Proceedings, Template, SBC OpenLib, Indexing
\end{keywords}

\begin{dates}
% This information will be provided by the editor before publishing the paper
\noindent{\sffamily\textbf{Recebido/Received:}} DD Month YYYY~~~$\bullet$~~~
{\sffamily\textbf{Aceito/Accepted:}} DD Month YYYY~~~$\bullet$~~~
{\sffamily\textbf{Publicado/Published:}} DD Month YYYY
\end{dates}


%\begin{license}
%Published under the Creative Commons Attribution 4.0 International Public License (CC BY 4.0)
%\end{license}

\end{frontmatter}


\section{Introdução}
\label{sec:intro}

O desenvolvimento de Veículos Aéreos Não Tripulados, popularmente conhecidos como \textit{drones}, tem crescido exponencialmente nas últimas décadas, impulsionado pelo avanço em áreas como eletrônica embarcada, sensores, comunicação sem fio e algoritmos de controle. Inicialmente restritos a aplicações militares, os VANTs tornaram-se elementos fundamentais em um vasto conjunto de contextos civis e científicos, como monitoramento ambiental, mapeamento geográfico, inspeções industriais, agricultura de precisão e resposta a desastres \cite{chung2018survey, hong2020swarm}. Essa expansão foi acompanhada por um aumento da autonomia e da capacidade de processamento dos sistemas embarcados, permitindo o surgimento de arquiteturas de controle mais sofisticadas e cooperativas.

Entretanto, apesar dos avanços, a operação isolada de um único VANT apresenta limitações inerentes, especialmente em termos de cobertura espacial, redundância e tolerância a falhas. Essas restrições motivaram a transição para abordagens cooperativas envolvendo múltiplos veículos, conhecidas como \textit{enxames de UAVs} (\textit{UAV swarms}) \cite{bu2024advancement}. Nesses sistemas, diversos VANTs atuam de maneira coordenada e distribuída, compartilhando informações e ajustando seus comportamentos locais de modo a atingir objetivos coletivos de forma mais eficiente e robusta do que um agente individual seria capaz de realizar.

O princípio fundamental dos enxames de UAVs inspira-se na \textit{inteligência coletiva} observada em sistemas naturais, como colônias de insetos, bandos de aves e cardumes de peixes. Essa abordagem, conhecida como \textit{Inteligência de Enxame} (\textit{Swarm Intelligence}), baseia-se em interações locais simples entre agentes, que, quando combinadas, dão origem a comportamentos globais emergentes \cite{ren2007consensus, chung2018survey}. Em um contexto de veículos aéreos, isso se traduz na capacidade de auto-organização, reconfiguração dinâmica e coordenação cooperativa, mesmo sob condições de comunicação limitada ou ambientes parcialmente desconhecidos.

Nesse cenário, o projeto de arquiteturas de controle adequadas é um dos principais desafios de pesquisa. Tais arquiteturas devem integrar comunicação eficiente, mecanismos de consenso, estratégias de controle de formação e coordenação de missão, equilibrando robustez, escalabilidade e custo computacional \cite{ouyang2023formation, bu2024advancement}. Além disso, a crescente demanda por sistemas multiagente em aplicações reais — como vigilância cooperativa, entrega autônoma e busca e salvamento — reforça a necessidade de modelos de simulação e validação que permitam avaliar o comportamento coletivo antes da implementação em campo.

Dessa forma, este trabalho tem como objetivo apresentar uma análise abrangente das principais abordagens, algoritmos e arquiteturas aplicadas ao controle cooperativo de enxames de VANTs. São discutidos conceitos fundamentais, como os níveis de autonomia e os princípios de sistemas multiagente, bem como metodologias de formação, estratégias de comunicação e desafios relacionados à coordenação distribuída. 

Com base nesses fundamentos, propõe-se investigar, por meio de simulações e testes comparativos, o desempenho de diferentes algoritmos de controle de formação — como \textit{Leader–Follower}, \textit{Consensus-Based}, \textit{Behavior-Based}, \textit{Virtual Structure} e \textit{Artificial Potential Fields (APF)} — em múltiplos cenários de operação. O estudo busca não apenas reproduzir resultados da literatura, mas avaliar criticamente a aplicabilidade e a eficiência de cada método sob condições diversas, incluindo ambientes com obstáculos, variações de topologia de comunicação e ajustes dinâmicos de parâmetros.

Por fim, também será explorada a possibilidade de hibridização entre algoritmos, de modo a combinar características complementares e identificar configurações capazes de otimizar a estabilidade da formação, a eficiência da comunicação e a adaptabilidade do enxame em missões cooperativas.



\section{Fundamentação Teórica}

O avanço tecnológico em eletrônicos, sensores e sistemas de comunicação tem impulsionado significativamente o desenvolvimento de Veículos Aéreos Não Tripulados (VANTs), popularmente conhecidos como drones, tornando-os cada vez mais acessíveis e aplicáveis em diversos contextos militares, civis e comerciais \cite{khan2017fanet}. Contudo, as limitações individuais de um único VANT — em termos de alcance, cobertura e carga útil — motivaram o surgimento de uma nova abordagem cooperativa: os enxames de VANTs (\textit{UAV swarms}) \cite{bu2024advancement, chung2018survey}.

Esses enxames consistem em múltiplos veículos coordenados de forma distribuída e descentralizada, capazes de executar tarefas coletivas de forma mais eficiente e robusta do que veículos isolados \cite{hong2020swarm}. Para que essa coordenação ocorra, é necessário o desenvolvimento de arquiteturas de controle robustas, capazes de integrar comunicação eficiente, formação dinâmica e cooperação autônoma entre os veículos \cite{hong2020swarm}.

\subsection{Conceitos Básicos}

\subsubsection{UAVs e Níveis de Autonomia}

Os VANTs têm apresentado rápida evolução nas últimas décadas, desempenhando funções em ambientes arriscados ou inacessíveis ao ser humano \cite{dong2004behavior}. Em muitas aplicações, um enxame de pequenos veículos pode apresentar desempenho superior ao de um único drone de grande porte, devido à sua flexibilidade, escalabilidade e redundância \cite{chung2018survey}.

A autonomia de um VANT refere-se à sua capacidade de operar de maneira independente, sem instruções humanas diretas \cite{sholes2007evolution}. Historicamente, os sistemas não tripulados se situavam entre dois extremos: \textit{teleoperação}, na qual o operador humano controla totalmente o veículo, e \textit{autonomia total}, em que o veículo toma decisões de forma independente \cite{dong2004behavior}.

Para quantificar esse grau de independência, \cite{sholes2007evolution} propôs uma taxonomia baseada nos níveis de controle de autonomia (ACL – Autonomy Control Level), estruturados segundo o ciclo Observar, Orientar, Decidir e Agir (OODA). Esse ciclo, originalmente concebido para o contexto militar, descreve a sequência de percepção, análise, tomada de decisão e ação de sistemas autônomos. Os níveis de ACL variam de 0, representando veículos totalmente teleoperados, a 10, correspondendo a sistemas plenamente autônomos e cientes de todo o ambiente de operação \cite{sholes2007evolution}.

O avanço da automação em VANTs implica a transferência de funções tradicionalmente desempenhadas por pilotos — como o planejamento de tarefas, a avaliação de eventos e a comunicação com estações terrestres — para o controlador de bordo \cite{dong2004behavior}.

\subsubsection{Sistemas Multiagente e Cooperação}

A coordenação entre múltiplos VANTs é geralmente tratada no contexto de Sistemas Multiagente (MAS – Multi-Agent Systems), nos quais diversos agentes autônomos interagem para atingir um objetivo comum \cite{bu2024advancement, chung2018survey}.

A cooperação é um princípio fundamental nesses sistemas. Quando bem coordenados, múltiplos veículos simples e de baixo custo podem alcançar desempenho global superior ao de um sistema monolítico \cite{ren2007consensus}. Em um MAS aplicado a VANTs, essa cooperação exige que funções como planejamento, adaptação a falhas e reação a emergências sejam parcialmente distribuídas entre os veículos \cite{dong2004behavior}.

\subsubsection{Inteligência de Enxame}

A Inteligência de Enxame (Swarm Intelligence) é inspirada em comportamentos coletivos observados na natureza, como em colônias de formigas, bandos de pássaros ou cardumes de peixes \cite{bu2024advancement, chung2018survey}.

O princípio central da inteligência de enxame está no comportamento emergente, que resulta das interações locais entre agentes simples, sem supervisão centralizada \cite{chung2018survey}. Essa abordagem distribuída é caracterizada pela troca local de informações entre vizinhos próximos (\textit{neighbor-to-neighbor information exchange}), permitindo que cada agente atualize seu estado com base nas informações recebidas e, coletivamente, o grupo alcance um consenso global \cite{ren2007consensus, bu2024advancement}.

\subsection{Comunicação e Protocolos entre UAVs}

A comunicação é um elemento essencial para a coordenação eficiente de enxames de VANTs. Para permitir a troca contínua de informações de estado, sensores e planos de voo, os VANTs frequentemente operam em redes ad hoc, conhecidas como FANETs (Flying Ad Hoc Networks) \cite{khan2017fanet, hong2020swarm}.

As FANETs diferenciam-se de redes terrestres tradicionais por características como alta mobilidade, auto-organização e ausência de controle centralizado. A dinâmica tridimensional dos VANTs, associada à sua velocidade de deslocamento (30 a 460 km/h), resulta em rápidas variações na topologia da rede, o que impõe desafios de conectividade e exige protocolos de comunicação resilientes \cite{khan2017fanet}.

Protocolos de comunicação comuns incluem Wi-Fi, UWB (Ultra-Wideband) e Bluetooth, selecionados conforme o alcance, largura de banda e consumo de energia requeridos \cite{chung2018survey}. A comunicação entre VANTs pode ser representada por meio de grafos direcionados ou não direcionados, em que cada nó representa um veículo e cada aresta representa um canal de comunicação ativo \cite{ren2007consensus}. Essa modelagem é particularmente útil para o desenvolvimento de algoritmos de consenso e controle distribuído.

\subsection{Arquiteturas de Controle para Enxames de UAVs}

A arquitetura de controle define como as decisões e informações são distribuídas entre os veículos do enxame. De modo geral, essas arquiteturas podem ser classificadas em três tipos principais: centralizada, descentralizada e distribuída \cite{bu2024advancement}.

\subsubsection{Controle Centralizado}

Na arquitetura centralizada, um controlador global — geralmente uma estação terrestre ou um VANT líder — é responsável por processar todas as informações e emitir comandos para os veículos subordinados \cite{ouyang2023formation, bu2024advancement}.

Essa abordagem é conceitualmente simples e permite decisões consistentes, além de facilitar a otimização global. Entretanto, sofre de limitações de escalabilidade e de robustez, já que a falha do controlador central pode comprometer toda a missão \cite{bu2024advancement}. Além disso, o processamento concentrado em um único ponto gera alta carga computacional, inviabilizando o uso em enxames muito grandes \cite{ouyang2023formation}.

\subsubsection{Controle Descentralizado e Distribuído}

No controle descentralizado, cada VANT toma decisões com base em informações locais, utilizando sensores embarcados ou comunicação com vizinhos próximos \cite{bu2024advancement}. Métodos como o \textit{Behavior-Based} e o \textit{Artificial Potential Field (APF)} são exemplos típicos dessa categoria.

Já no controle distribuído, embora não exista um centro de comando, há cooperação explícita entre os VANTs, que compartilham informações de estado para alcançar objetivos coletivos. Essa abordagem aumenta a robustez do sistema e melhora a escalabilidade, mas exige algoritmos mais complexos de consenso e sincronização \cite{bu2024advancement, ouyang2023formation}.

\subsubsection{Comparação entre Arquiteturas}

A Tabela~\ref{tab:arquiteturas-uavs} resume as principais características das arquiteturas de controle apresentadas anteriormente, destacando diferenças em robustez, escalabilidade e complexidade de implementação. Observa-se que a arquitetura centralizada apresenta baixa tolerância a falhas e escalabilidade limitada, porém sua implementação é conceitualmente mais simples. Por outro lado, arquiteturas descentralizadas ou distribuídas oferecem maior robustez e escalabilidade, mas exigem algoritmos de coordenação mais complexos, refletindo o trade-off clássico entre simplicidade de design e desempenho do sistema.


\begin{table*}[h!]
\centering
\caption{Comparação entre arquiteturas de controle em enxames de UAVs.}
\label{tab:arquiteturas-uavs}
\begin{tabular}{|p{3cm}|p{4cm}|p{3cm}|p{4cm}|}
\hline
\textbf{Estratégia} & \textbf{Robustez} & \textbf{Escalabilidade} & \textbf{Complexidade} \\ \hline
\textbf{Centralizada} & Baixa (alto risco de ponto único de falha) \cite{ouyang2023formation} & Limitada (sobrecarga computacional) \cite{bu2024advancement} & Simples (decisão unificada) \cite{bu2024advancement} \\ \hline
\textbf{Descentralizada e Distribuída} & Alta (forte tolerância a falhas) \cite{ouyang2023formation} & Alta (baseada em interações locais) \cite{bu2024advancement} & Elevada (coordenação e comunicação complexas) \cite{bu2024advancement} \\ \hline
\end{tabular}
\end{table*}


O controle distribuído tem se destacado como uma das tendências mais promissoras na pesquisa contemporânea, devido à sua robustez, flexibilidade e capacidade de escala \cite{ouyang2023formation}. Essa abordagem elimina a dependência de um ponto central e permite a coordenação emergente entre os veículos, sendo considerada a direção mais provável para o futuro das missões cooperativas com enxames de UAVs.

\subsection{Algoritmos de Formação de Drones}

O controle de formação de \textit{Unmanned Aerial Vehicles} (UAVs) consiste em coordenar o movimento de múltiplos veículos de modo que mantenham uma disposição espacial específica, geralmente representada por padrões geométricos como linha, V, triângulo ou diamante \cite{bu2024advancement}. Essa tarefa é central em aplicações cooperativas, como vigilância, mapeamento, busca e salvamento, e monitoramento ambiental \cite{bu2024advancement, ouyang2023formation}.

Matematicamente, o problema de formação pode ser modelado considerando um conjunto de $N$ drones com posição $p_i \in \mathbb{R}^n$ e velocidade $v_i \in \mathbb{R}^n$, onde $i = 1, \dots, N$. O objetivo é garantir que cada par de drones $(i, j)$ mantenha a distância e orientação relativas desejadas, ou seja:

\[
| (p_i - p_j) - d_{ij} | \to 0, \quad \forall (i, j) \in E
\]

onde $d_{ij}$ representa o vetor desejado de separação entre os drones $i$ e $j$, e $E$ é o conjunto de conexões (vizinhanças) entre os UAVs \cite{bu2024advancement, ouyang2023formation}.

Diversas abordagens foram propostas na literatura para alcançar o controle de formação. Entre as mais estudadas estão: \textit{Leader–Follower}, \textit{Virtual Structure}, \textit{Behavior-Based}, \textit{Consensus-Based} e \textit{Artificial Potential Fields}. Cada uma apresenta vantagens e limitações específicas quanto à escalabilidade, robustez, custo computacional e dependência de comunicação entre os veículos \cite{bu2024advancement, ouyang2023formation}.

\subsubsection{Leader–Follower}

A abordagem \textit{Leader–Follower} baseia-se em uma hierarquia onde um ou mais UAVs assumem o papel de líderes, enquanto os demais atuam como seguidores, mantendo posições relativas fixas em relação aos líderes \cite{bu2024advancement}. Essa técnica é intuitiva e de fácil implementação, sendo amplamente utilizada em formações com poucos agentes.

O erro de posição para um seguidor $i$ em relação a um líder $L$ é dado por:

\[
e_i = (p_L + R_L d_i) - p_i
\]

onde $R_L$ representa a orientação do líder e $d_i$ o vetor desejado da posição relativa do seguidor. O controle de velocidade é definido como:

\[
\dot{p}_i = v_i = \dot{p}_L + K_p e_i
\]

com $K_p > 0$ representando o ganho proporcional que determina a taxa de correção do erro \cite{bu2024advancement}. Essa abordagem, entretanto, é sensível à perda de comunicação com o líder, podendo comprometer toda a formação \cite{ouyang2023formation}.


\subsubsection{Leader–Follower}

A abordagem \textit{Leader–Follower} baseia-se em uma estrutura hierárquica em que um ou mais UAVs atuam como líderes, definindo a trajetória global, enquanto os demais agentes (seguidores) mantêm posições relativas fixas em relação a eles. Essa técnica é amplamente empregada por sua simplicidade e eficiência em pequenas formações.

Matematicamente, a posição desejada de um seguidor $i$ em relação a um líder $L$ é expressa por

\[
p_i^{d}(t) = p_L(t) + R_L(t) \, d_i,
\]

onde $p_L(t)$ representa a posição do líder, $R_L(t) \in SO(3)$ é a matriz de rotação correspondente à sua orientação, e $d_i$ é o vetor fixo da posição desejada do seguidor no referencial do líder. A inclusão de $R_L(t)$ permite que a formação mantenha sua geometria mesmo quando o líder altera sua atitude.

O erro de formação é definido como

\[
e_i(t) = p_i^{d}(t) - p_i(t) = p_L(t) + R_L(t)d_i - p_i(t),
\]

e a lei de controle de velocidade, considerando um esquema proporcional-derivativo (PD) com \textit{feed-forward} da velocidade do líder, é dada por

\[
u_i(t) = v_L(t) + K_p e_i(t) + K_v \big(v_L(t) - v_i(t)\big),
\]

em que $u_i(t)$ é o comando de velocidade do seguidor, $v_L(t) = \dot{p}_L(t)$ é a velocidade do líder, e $K_p, K_v > 0$ são ganhos de controle que ajustam a resposta dinâmica da formação. Essa formulação assegura que os seguidores mantenham a distância e orientação desejadas em relação ao líder, acompanhando suas variações de trajetória de forma suave e estável.

Apesar de sua simplicidade, o método é sensível a falhas de comunicação ou perda do líder, o que pode comprometer a coesão do grupo. Estratégias recentes mitigam essa limitação com múltiplos líderes ou controle cooperativo distribuído.


\subsubsection{Virtual Structure}

Na estratégia \textit{Virtual Structure}, o grupo de UAVs é tratado como uma estrutura rígida virtual, na qual cada veículo é responsável por ocupar uma posição fixa dentro dessa estrutura \cite{bu2024advancement}. Dessa forma, o enxame se move como uma entidade unificada, garantindo precisão na manutenção da formação.

Seja $p_v(t)$ a posição da estrutura virtual e $R_v(t)$ sua orientação. A posição desejada do drone $i$ é dada por:

\[
p_i^d(t) = p_v(t) + R_v(t)d_i
\]

O controle de posição é então expresso como:

\[
\dot{p}_i = \dot{p}_v + R_v \dot{d}_i + K_p (p_i^d - p_i)
\]

Essa técnica proporciona excelente estabilidade e coerência na formação, porém apresenta baixa adaptabilidade a ambientes com obstáculos ou perturbações externas \cite{bu2024advancement, ouyang2023formation}.


\subsubsection{Virtual Structure}

Na estratégia \textit{Virtual Structure}, o grupo de UAVs é modelado como uma única estrutura rígida virtual, na qual cada veículo ocupa uma posição fixa pré-definida dentro dessa estrutura \cite{bu2024advancement}. Assim, o enxame se desloca como uma entidade unificada, mantendo a forma e a orientação da formação com elevada precisão.

Seja $p_v(t) \in \mathbb{R}^3$ a posição do centro da estrutura virtual e $R_v(t) \in SO(3)$ sua matriz de rotação, que define a orientação global da formação. Cada UAV $i$ está associado a um vetor $r_i^{VS}$, que representa sua posição relativa fixa em relação ao centro da estrutura. Dessa forma, a posição desejada do UAV $i$ é dada por:

\[
p_i^{d}(t) = p_v(t) + R_v(t) r_i^{VS}
\]

A velocidade desejada correspondente é obtida pela cinemática de corpo rígido:

\[
v_i^{d}(t) = v_v(t) + \omega_v(t) \times R_v(t) r_i^{VS}
\]

onde $v_v(t)$ é a velocidade linear e $\omega_v(t)$ é a velocidade angular da estrutura virtual. 

O controle de posição e velocidade de cada UAV é definido por uma lei de controle proporcional-derivativa (PD), dada por:

\[
u_i = K_p (p_i^{d} - p_i) + K_v (v_i^{d} - v_i)
\]

onde $u_i$ representa o vetor de aceleração de controle do UAV $i$, e $K_p$, $K_v$ são matrizes de ganhos proporcionais e derivativos, respectivamente. Substituindo $v_i^{d}(t)$, tem-se explicitamente:

\[
u_i = K_p (p_v + R_v r_i^{VS} - p_i) + K_v (v_v + \omega_v \times R_v r_i^{VS} - v_i)
\]

Essa lei de controle assegura que cada UAV siga rigidamente o movimento da estrutura virtual, compensando desvios de posição e velocidade em relação à sua referência local dentro do corpo virtual.

Assumindo uma dinâmica de segunda ordem $\ddot{p}_i = u_i$, o erro de posição $e_i = p_i^{d} - p_i$ evolui segundo:

\[
\ddot{e}_i + K_v \dot{e}_i + K_p e_i = 0
\]

o que caracteriza um sistema linear de segunda ordem estável se $K_p > 0$ e $K_v > 0$, garantindo convergência exponencial do erro para zero.

A principal vantagem dessa abordagem é a manutenção precisa da geometria da formação, mesmo em manobras complexas. Entretanto, ela requer sincronização e comunicação global entre os UAVs para garantir consistência no estado da estrutura virtual, o que pode limitar sua aplicação em cenários com comunicação intermitente ou ambientes altamente dinâmicos \cite{bu2024advancement, ouyang2023formation}.




\subsubsection{Behavior-Based}

A abordagem \textit{Behavior-Based} fundamenta-se em regras locais simples, inspiradas no comportamento coletivo de animais, como aves ou peixes. Cada UAV reage a seus vizinhos próximos de acordo com princípios de \textit{separação}, \textit{coesão} e \textit{alinhamento}, levando a um comportamento global emergente \cite{ouyang2023formation}.

A força de controle total aplicada ao drone $i$ é uma soma ponderada de cada comportamento:

\[
u_i = w_s f_{sep,i} + w_c f_{coh,i} + w_a f_{ali,i}
\]

onde:

\[
\begin{aligned}
f_{sep,i} &= \sum_{j \in N_i} \frac{p_i - p_j}{|p_i - p_j|^2} \quad &\text{(Separação)} \\[6pt]
f_{coh,i} &= \frac{1}{|N_i|}\sum_{j \in N_i}(p_j - p_i) \quad &\text{(Coesão)} \\[6pt]
f_{ali,i} &= \frac{1}{|N_i|}\sum_{j \in N_i}(v_j - v_i) \quad &\text{(Alinhamento)}
\end{aligned}
\]

O movimento é então atualizado por:

\[
\dot{p}_i = v_i, \quad \dot{v}_i = u_i
\]

Essa abordagem é descentralizada e robusta a falhas de comunicação, mas tende a ter dificuldades em alcançar formações geométricas rígidas e precisas \cite{bu2024advancement, ouyang2023formation}.


\subsubsection{Behavior-Based}

A abordagem \textit{Behavior-Based} fundamenta-se em regras locais simples, inspiradas em sistemas biológicos como bandos de aves e cardumes de peixes. Cada UAV ajusta seu movimento de forma autônoma com base em informações provenientes apenas de seus vizinhos mais próximos, seguindo comportamentos fundamentais de \textit{separação}, \textit{coesão}, \textit{alinhamento} e, em algumas aplicações, \textit{busca por objetivo} (\textit{goal seeking}) \cite{ouyang2023formation, reynolds1987flocks}.

A força de controle total aplicada ao agente $i$ é obtida pela soma ponderada dos vetores de cada comportamento:

\[
u_i = w_s f_{sep,i} + w_c f_{coh,i} + w_a f_{ali,i} + w_g f_{goal,i}
\]

onde $w_s$, $w_c$, $w_a$ e $w_g$ são os pesos que determinam a influência relativa de cada comportamento. Os vetores de controle são definidos como:

\[
\begin{aligned}
f_{sep,i} &= \sum_{j \in N_i} \frac{p_i - p_j}{\|p_i - p_j\|^2} && \text{(Separação)} \\[6pt]
f_{coh,i} &= \frac{1}{|N_i|}\sum_{j \in N_i}(p_j - p_i) && \text{(Coesão)} \\[6pt]
f_{ali,i} &= \frac{1}{|N_i|}\sum_{j \in N_i}(v_j - v_i) && \text{(Alinhamento)} \\[6pt]
f_{goal,i} &= p_{goal} - p_i && \text{(Busca por objetivo)}
\end{aligned}
\]

O comportamento de \textit{separação} evita colisões e mantém espaçamento mínimo seguro entre os UAVs; o de \textit{coesão} busca preservar a unidade do grupo; o de \textit{alinhamento} sincroniza a direção de voo com a média dos vizinhos; e o de \textit{goal seeking} conduz o enxame a um alvo ou trajetória desejada.

A dinâmica de cada agente é descrita por:

\[
\dot{p}_i = v_i, \quad \dot{v}_i = u_i
\]

onde $p_i$ e $v_i$ são, respectivamente, a posição e a velocidade do UAV $i$. Assim, o vetor $u_i$ representa a aceleração desejada calculada a partir da combinação dos comportamentos locais.

Essa formulação resulta em um sistema de controle totalmente distribuído, no qual o comportamento coletivo emerge das interações locais entre agentes. Tal característica garante escalabilidade e robustez a falhas de comunicação, mas dificulta a manutenção de formações geométricas rígidas e a obtenção de estabilidade formal em cenários complexos \cite{bu2024advancement, ouyang2023formation}.



\subsubsection{Consensus-Based}

Os métodos \textit{Consensus-Based} visam permitir que todos os UAVs atinjam um consenso sobre variáveis de estado, como posição ou velocidade, usando apenas informações locais de vizinhança \cite{bu2024advancement, ouyang2023formation}. A dinâmica contínua de consenso é expressa como:

\[
\dot{x}_i = -\sum_{j \in N_i} a_{ij}(x_i - x_j)
\]

onde $x_i$ é o estado do agente $i$ e $a_{ij}$ é o peso da comunicação entre os agentes. Em forma matricial, o sistema é representado por:

\[
\dot{x} = -Lx
\]

sendo $L$ a matriz Laplaciana do grafo de comunicação. Quando o grafo é fortemente conexo, o sistema garante:

\[
\lim_{t \to \infty} x_i(t) = x_j(t), \quad \forall i, j
\]

ou seja, todos os drones convergem para o mesmo valor \cite{ouyang2023formation}. Essa técnica possui base teórica sólida e é amplamente utilizada em sistemas distribuídos pela sua robustez e simplicidade \cite{bu2024advancement}.

\subsubsection{Consensus-Based}

Os métodos \textit{Consensus-Based} têm como objetivo permitir que todos os UAVs atinjam um acordo (\textit{consenso}) sobre variáveis de estado, como posição, velocidade ou orientação, utilizando apenas informações locais de vizinhança \cite{bu2024advancement, ouyang2023formation, olfati2007consensus, ren2008distributed}. 

A topologia de comunicação entre os UAVs é representada por um grafo direcionado $\mathcal{G} = (\mathcal{V}, \mathcal{E})$, onde $\mathcal{V} = \{1, 2, \dots, N\}$ é o conjunto de agentes e $\mathcal{E} \subseteq \mathcal{V} \times \mathcal{V}$ define as conexões entre eles. A matriz de adjacência $A = [a_{ij}]$ é composta por pesos de comunicação, sendo $a_{ij} > 0$ se o agente $j$ envia informações para o agente $i$, e $a_{ij} = 0$ caso contrário. A matriz Laplaciana é definida como $L = D - A$, em que $D = \text{diag}(d_1, d_2, \dots, d_N)$, com $d_i = \sum_{j} a_{ij}$.

\paragraph{Modelagem para agentes de primeira ordem.}
Considerando que cada agente $i$ possui um estado $x_i(t)$ e que sua dinâmica é dada por $\dot{x}_i = u_i$, o controle de consenso é expresso por:
\[
u_i = -\sum_{j \in \mathcal{N}_i} a_{ij} (x_i - x_j)
\]
Substituindo em $\dot{x}_i = u_i$, tem-se a forma dinâmica contínua:
\[
\dot{x}_i = -\sum_{j \in \mathcal{N}_i} a_{ij} (x_i - x_j)
\]
ou, de forma vetorial compacta:
\[
\dot{x} = -Lx
\]
onde $x = [x_1, x_2, \dots, x_N]^T$. Quando o grafo é fortemente conexo, o sistema garante a convergência para o consenso global:
\[
\lim_{t \to \infty} (x_i(t) - x_j(t)) = 0, \quad \forall i,j
\]

\paragraph{Modelagem para agentes de segunda ordem.}
Em sistemas mais realistas, como enxames de UAVs, é necessário considerar tanto a posição $p_i$ quanto a velocidade $v_i$ de cada agente. Assim, a dinâmica de segunda ordem é dada por:
\[
\begin{cases}
\dot{p}_i = v_i \\
\dot{v}_i = u_i
\end{cases}
\]
O controle de consenso estendido para esse caso pode ser expresso como:
\[
u_i = -k_p \sum_{j \in \mathcal{N}_i} a_{ij}(p_i - p_j) - k_v \sum_{j \in \mathcal{N}_i} a_{ij}(v_i - v_j)
\]
em que $k_p$ e $k_v$ são ganhos de acoplamento que determinam a velocidade de convergência da formação. Sob condições de conectividade adequadas do grafo, o sistema tende a garantir:
\[
\lim_{t \to \infty} (p_i(t) - p_j(t)) = 0, \quad \lim_{t \to \infty} (v_i(t) - v_j(t)) = 0, \quad \forall i,j
\]
ou seja, todos os UAVs convergem para a mesma trajetória e velocidade, atingindo consenso tanto em posição quanto em movimento.

\paragraph{Discussão.}
Essa abordagem apresenta uma base teórica sólida e comprovada estabilidade quando o grafo de comunicação é conexo. Sua simplicidade, robustez e escalabilidade tornam o método \textit{Consensus-Based} amplamente utilizado em aplicações de coordenação distribuída, controle de formação e sincronização em redes multiagente \cite{olfati2007consensus, ren2008distributed}.



\subsubsection{Artificial Potential Fields}

O método dos \textit{Artificial Potential Fields} (APF) propõe o uso de campos de potencial artificiais para gerar forças virtuais de atração e repulsão entre os agentes \cite{bu2024advancement, ouyang2023formation}. O potencial total de um drone $i$ é definido como:

\[
U_i = U_{att}(p_i) + U_{rep}(p_i)
\]

com:

\[
U_{att}(p_i) = \frac{1}{2} k_{att} |p_i - p_{goal}|^2
\]

\[
U_{rep}(p_i) = \frac{1}{2} k_{rep} \sum_{j \in N_i} \left( \frac{1}{|p_i - p_j|} - \frac{1}{d_0} \right)^2 \mathbb{I}_{|p_i - p_j| < d_0}
\]

A força resultante que guia o movimento é:

\[
f_i = -\nabla U_i = f_{att,i} + f_{rep,i}
\]

Esse método é intuitivo e eficaz para evitar colisões e reorganizar formações em tempo real, mas pode sofrer com o problema dos mínimos locais, onde os UAVs ficam presos em posições de equilíbrio indesejadas \cite{bu2024advancement, ouyang2023formation}.


\subsubsection{Artificial Potential Fields}

O método dos \textit{Artificial Potential Fields} (APF) propõe o uso de campos de potencial artificiais para gerar forças virtuais de atração e repulsão que guiam o movimento dos agentes de forma contínua e descentralizada \cite{khatib1986real, bu2024advancement, ouyang2023formation}. Nesse contexto, o ambiente é modelado como um campo energético, em que o objetivo atua como uma fonte atrativa e os obstáculos ou demais UAVs funcionam como fontes repulsivas.

O potencial total associado ao agente $i$ é definido por:

\[
U_i(p_i) = U_{att}(p_i) + U_{rep}(p_i)
\]

em que o termo atrativo direciona o UAV para o alvo e o termo repulsivo evita colisões com obstáculos e vizinhos próximos. O potencial atrativo é geralmente modelado de forma quadrática:

\[
U_{att}(p_i) = \frac{1}{2} k_{att} \|p_i - p_{goal}\|^2
\]

onde $k_{att} > 0$ representa o ganho de atração e $p_{goal}$ a posição desejada. Já o potencial repulsivo atua apenas dentro de uma distância de influência $d_0$, sendo dado por:

\[
U_{rep}(p_i) =
\frac{1}{2} k_{rep} 
\sum_{j \in N_i}
\left(
\frac{1}{\|p_i - p_j\|} - \frac{1}{d_0}
\right)^2
\mathbb{I}_{\|p_i - p_j\| < d_0}
\]

onde $k_{rep} > 0$ é o ganho de repulsão, $N_i$ representa o conjunto de vizinhos do agente $i$, e $\mathbb{I}$ é a função indicadora, que vale 1 quando a condição é verdadeira e 0 caso contrário.

A força resultante que governa o movimento do agente é obtida pelo gradiente negativo do potencial total:

\[
f_i = -\nabla U_i = f_{att,i} + f_{rep,i}
\]

com:

\begin{equation}
\begin{split}
f_{att,i} &= -k_{att}(p_i - p_{goal}), \\
f_{rep,i} &= 
\sum_{j \in N_i} k_{rep}
\left(\frac{1}{\|p_i - p_j\|} - \frac{1}{d_0}\right)
\\
&\quad\times
\frac{1}{\|p_i - p_j\|^3}(p_i - p_j)
\mathbb{I}_{\|p_i - p_j\| < d_0}
\end{split}
\end{equation}

A dinâmica de movimento pode ser modelada por:

\[
\dot{p}_i = v_i, \quad
\dot{v}_i = f_i
\]

de modo que cada UAV ajusta sua aceleração de acordo com o gradiente do campo de potencial. Essa formulação permite que as formações se reorganizem em tempo real e evitem colisões, mesmo em ambientes dinâmicos. No entanto, o método pode apresentar limitações devido à presença de mínimos locais — situações em que o potencial total apresenta pontos de equilíbrio indesejados que impedem a convergência ao objetivo \cite{khatib1986real, bu2024advancement, ouyang2023formation}.

\subsubsection*{Comparação Geral}

Cada uma dessas abordagens apresenta vantagens específicas. Os métodos \textit{Leader–Follower} e \textit{Virtual Structure} são indicados para formações bem definidas e controladas centralmente, enquanto as estratégias \textit{Behavior-Based}, \textit{Consensus-Based} e \textit{Artificial Potential Fields} são mais adequadas para ambientes dinâmicos e distribuídos. Em sistemas reais, é comum a combinação híbrida entre métodos, buscando equilibrar estabilidade, escalabilidade e adaptabilidade \cite{bu2024advancement, ouyang2023formation}.

\subsection{Comunicação e Coordenação em Enxames de UAVs}

A comunicação e a coordenação são pilares fundamentais para o funcionamento de sistemas cooperativos de VANTs (Veículos Aéreos Não Tripulados). Esses elementos determinam não apenas a eficiência das formações, mas também a robustez e escalabilidade das missões em ambientes dinâmicos e distribuídos \cite{ren2007consensus, Young-Min2018, khaliq2021}.

Através de protocolos de comunicação específicos, como o \textit{MAVLink}, e estruturas de software distribuídas, como o \textit{ROS} e o \textit{ROS2}, é possível garantir a troca contínua de dados entre os veículos e entre estes e a estação de controle \cite{ziming2023}. No entanto, para que a cooperação ocorra de maneira eficiente, é essencial que a rede de comunicação suporte o compartilhamento rápido e confiável de informações — algo desafiador em redes altamente dinâmicas, como as \textit{FANETs} (Flying Ad hoc Networks), sujeitas à latência, perdas de pacotes e instabilidade topológica \cite{inam2017, hong2020swarm, ouyang2023formation}.

\subsubsection{Protocolos UAV e Arquiteturas de Comunicação}

O protocolo \textit{MAVLink}, amplamente utilizado em sistemas \textit{open-source} como PX4 e ArduPilot, provê uma camada leve de comunicação baseada em mensagens que permite a troca de dados entre o sistema de controle de voo e a estação terrestre \cite{Young-Min2018}. Em conjunto, o \textit{ROS} e o \textit{ROS2} fornecem uma infraestrutura modular e distribuída que viabiliza a integração de múltiplos agentes, sensores e atuadores, promovendo interoperabilidade entre diferentes sistemas e linguagens \cite{khaliq2021, ziming2023}.

Essas tecnologias permitem o desenvolvimento de arquiteturas descentralizadas, nas quais cada UAV pode atuar como um nó inteligente em uma rede distribuída, trocando informações localmente com seus vizinhos e ajustando seu comportamento coletivo com base em protocolos de consenso \cite{ren2007consensus, bu2024advancement}. O uso do \textit{DDS} (Data Distribution Service) em \textit{ROS2} melhora a escalabilidade e a confiabilidade dessas trocas, tornando o sistema mais tolerante a falhas de comunicação.

\subsubsection{Sincronização e Coordenação de Estados}

A sincronização entre múltiplos UAVs é obtida por meio de algoritmos de controle cooperativo, especialmente aqueles baseados em consenso \cite{ren2007consensus}. Esses métodos asseguram que os agentes compartilhem informações locais suficientes para alcançar um estado comum — como alinhamento de velocidade, posição ou orientação — sem necessidade de um controlador centralizado \cite{bu2024advancement}.

Em essência, o consenso é o mecanismo pelo qual um enxame “pensa junto”. Cada drone ajusta suas variáveis de estado em função dos vizinhos, e, sob certas condições de conectividade do grafo de comunicação, o sistema converge para um comportamento coletivo estável e coerente \cite{chung2018survey}. Essa propriedade é explorada em diversos paradigmas de formação, como o modelo \textit{leader–follower}, o de \textit{estrutura virtual} e o de \textit{behavior-based}, todos fundamentados em princípios semelhantes de cooperação e sincronização distribuída.

\subsubsection{Limitações e Desafios de Comunicação}

Apesar do avanço nos protocolos e arquiteturas, a comunicação entre UAVs permanece sujeita a limitações significativas. A natureza tridimensional e dinâmica das \textit{FANETs} impõe desafios quanto à latência, perda de pacotes e vulnerabilidade a ataques de rede, como \textit{ICMP flooding} e negação de serviço (DoS) \cite{Young-Min2018}. Esses fatores podem comprometer a estabilidade de formações e causar falhas críticas em missões coordenadas.

Estudos experimentais mostraram que ataques de saturação podem aumentar a variância do tempo de recepção de dados em mais de 30 vezes, afetando diretamente o desempenho dos controladores de formação \cite{Young-Min2018}. Além disso, a degradação do sinal (RSSI) e a quebra de links em redes \textit{mesh} reduzem a capacidade de sincronização entre os drones, exigindo estratégias adaptativas de roteamento e controle \cite{hong2020swarm, khaliq2021}.

\subsection{Simulação de Enxames de UAVs}

As plataformas de simulação são ferramentas essenciais no desenvolvimento e validação de estratégias de controle cooperativo e comunicação entre UAVs \cite{khaliq2021, ziming2023}. Elas permitem testar algoritmos de formação, sincronização e evasão de obstáculos em ambientes virtuais, reduzindo custos e riscos associados a testes reais.

\subsubsection{Plataformas e Motores de Simulação}

Simuladores como \textit{Gazebo}, \textit{Webots}, \textit{CoppeliaSim} e o \textit{SITL PX4} fornecem ambientes tridimensionais realistas para a modelagem de cenários de voo e formação \cite{ziming2023}. O \textit{Gazebo}, por exemplo, integra-se ao \textit{ROS} para simulações \textit{Software-In-The-Loop} (SIL) e \textit{Hardware-In-The-Loop} (HIL), utilizando motores de física como ODE, Bullet e DART \cite{khaliq2021}.

Já o \textit{Webots} destaca-se pela flexibilidade e suporte a sensores complexos — como LIDAR, GPS e IMU — sendo amplamente utilizado em experimentos de enxames e controle distribuído \cite{ziming2023}. Tais simuladores não apenas reproduzem o comportamento físico dos UAVs, mas também emulam suas redes de comunicação via MAVLink e ROS, permitindo testes completos de sistemas multiagentes.

\subsubsection{Comunicação e Modelagem em Simulação}

Durante simulações cooperativas, o MAVLink atua como o protocolo de mensageria entre os nós simulados, e o MAVROS realiza a interface entre o sistema de controle e os tópicos ROS \cite{khaliq2021}. Essa configuração possibilita a criação de redes \textit{mesh} distribuídas, nas quais cada UAV emulado executa seu próprio nó ROS e interage com os demais por meio de mensagens trocadas em tempo real.

Além disso, a modelagem de sensores e ambientes permite testar a robustez dos algoritmos em condições realistas, incluindo ruído, atraso de comunicação e falhas de hardware simuladas \cite{ziming2023}. Essa integração entre simulação física e rede distribuída fornece uma base sólida para o desenvolvimento de estratégias de coordenação robustas.

\subsection{Síntese e Direcionamento para o Trabalho}

A fundamentação teórica apresentada integrou os principais conceitos relacionados à comunicação, coordenação, simulação e controle de formação em enxames de UAVs, delineando o panorama atual das tecnologias e desafios envolvidos no desenvolvimento de sistemas cooperativos autônomos.

Os protocolos de comunicação, como o \textit{MAVLink} e o \textit{ROS/ROS2}, fornecem a base para a interoperabilidade e sincronização entre múltiplos agentes, permitindo a troca de informações de estado e o controle distribuído. Sobre essa infraestrutura, os algoritmos de formação — como os baseados em consenso, comportamento e estruturas virtuais — determinam a dinâmica coletiva dos UAVs, regulando como o grupo se organiza, mantém coesão e responde a perturbações externas. Já as plataformas de simulação, como o \textit{Gazebo} e o \textit{SITL PX4}, tornam possível testar e validar essas estratégias em cenários controlados, antes de sua aplicação em experimentos reais \cite{ouyang2023formation, bu2024advancement, khaliq2021}.

Com base nessa revisão, este trabalho propõe investigar de forma comparativa o desempenho de diferentes algoritmos e arquiteturas de controle de formação em múltiplos contextos experimentais. O foco não se limita à reprodução de resultados consolidados, mas à análise crítica da aplicabilidade de cada método sob condições diversas — como a presença de obstáculos, variações de topologia de comunicação e ajustes de parâmetros dinâmicos.

Além disso, será explorada a hibridização de algoritmos, combinando propriedades de estratégias como \textit{Leader–Follower}, \textit{Consensus-Based} e \textit{Artificial Potential Fields (APF)}. Essa abordagem busca identificar configurações que maximizem a estabilidade da formação e a eficiência da comunicação, sem comprometer a adaptabilidade e a robustez do enxame.

\section{Metodologia}

A metodologia deste trabalho foi desenvolvida com o objetivo de realizar uma análise comparativa entre diferentes algoritmos de controle cooperativo de formação de VANTs (Veículos Aéreos Não Tripulados). A pesquisa adota uma abordagem experimental e quantitativa, baseada em simulações computacionais realizadas em ambientes virtuais. O processo geral é dividido em quatro etapas principais: definição dos algoritmos e cenários, implementação dos modelos de controle, execução das simulações e análise dos resultados.

Inicialmente, são definidos os algoritmos que compõem o escopo do estudo, abrangendo as principais abordagens de controle cooperativo de formação:

\begin{itemize}
    \item \textbf{Leader–Follower}: um agente líder define a trajetória e os seguidores ajustam suas posições relativas com base nele;
    \item \textbf{Consensus-Based}: todos os agentes atualizam seus estados (posição, velocidade) de acordo com os vizinhos, buscando convergência coletiva;
    \item \textbf{Behavior-Based}: o comportamento global emerge da combinação de ações individuais, como coesão, separação e alinhamento;
    \item \textbf{Virtual Structure}: a formação é tratada como uma estrutura rígida virtual, na qual cada agente segue uma posição relativa predefinida;
    \item \textbf{Artificial Potential Fields (APF)}: os drones se movem segundo forças atrativas e repulsivas que determinam o posicionamento desejado.
\end{itemize}

Além da análise individual, também serão avaliadas combinações híbridas de algoritmos, explorando a complementaridade entre estratégias — por exemplo, a integração de \textit{Consensus-Based} com \textit{APF} para manter a coesão da formação e, simultaneamente, evitar colisões.

As simulações serão realizadas em ambiente virtual utilizando o simulador \textit{Gazebo} integrado ao middleware \textit{ROS2}, com o firmware PX4 Autopilot em modo SITL (\textit{Software In The Loop}) e o protocolo MAVLink para comunicação entre os agentes. As implementações serão desenvolvidas em \textit{Python}, com o uso de bibliotecas como \textit{NumPy} e \textit{Matplotlib} para tratamento e visualização dos dados.

Três cenários principais serão configurados:
\begin{enumerate}
    \item \textbf{Espaço Livre:} ambiente sem obstáculos, para avaliação da formação e manutenção da estrutura;
    \item \textbf{Com Obstáculos Estáticos:} presença de barreiras fixas, testando a reconfiguração da formação;
    \item \textbf{Com Obstáculos Dinâmicos:} inclusão de obstáculos móveis, avaliando a capacidade de resposta e adaptação do enxame.
\end{enumerate}

Cada cenário será simulado múltiplas vezes com condições iniciais distintas, assegurando robustez estatística. As métricas de desempenho consideradas incluem:

\begin{itemize}
    \item \textbf{Erro de Formação (EF):} distância média entre a posição real e a posição desejada de cada agente;
    \item \textbf{Tempo de Convergência (TC):} tempo necessário para o sistema atingir uma configuração estável;
    \item \textbf{Estabilidade e Robustez (ER):} capacidade do sistema de manter a formação diante de falhas de comunicação ou distúrbios externos.
\end{itemize}

A análise dos resultados será conduzida de forma quantitativa e gráfica, por meio de curvas de convergência, gráficos de dispersão e estatísticas descritivas (média, variância e desvio padrão) para cada métrica. Essa comparação permitirá identificar quais algoritmos — ou combinações híbridas — apresentam melhor desempenho em diferentes contextos.

Por fim, a metodologia proposta busca garantir reprodutibilidade e comparabilidade entre os experimentos. A partir dela, pretende-se identificar estratégias cooperativas mais estáveis, eficientes e adaptáveis, contribuindo para o avanço do controle distribuído em enxames de UAVs.


\section{Conclusão}

\begin{declarations}

\begin{acknowledgements}
ESTA DECLARAÇÃO É OPCIONAL. Este é um texto de agradecimentos com várias linhas. Lorem ipsum dolor sit amet, consectetur adipiscing elit, sed do eiusmod tempor incididunt ut labore et dolore magna aliqua. Ut enim ad minim veniam, quis nostrud exercitation ullamco laboris nisi ut aliquip ex ea commodo consequat.
\end{acknowledgements}


\begin{funding}
ESTA DECLARAÇÃO É OPCIONAL. Esta pesquisa foi financiada por lorem ipsum dolor sit amet, consectetur adipiscing elit.
\end{funding}

\begin{contributions}
ESTA DECLARAÇÃO É OBRIGATÓRIA. 
\textit{Considerar o Código de Conduta da SBC.
Participação em autoria: espera-se que todos os autores de um trabalho publicado 
ou submetido para publicação, tenham tido efetiva participação no respectivo trabalho.}
Sugerimos que os autores descrevam sua contribuição usando 
a Taxonomia CRediT (\href{https://credit.niso.org/}{https://credit.niso.org/}) como neste exemplo: 
JV contribuiu para a concepção deste estudo. CB, RP e CM realizaram os experimentos. 
JV é o principal contribuidor e escritor deste manuscrito. 
Todos os autores leram e aprovaram o manuscrito final. 
\end{contributions}

\begin{interests}
ESTA DECLARAÇÃO É OBRIGATÓRIA. 
Se não houver conflitos de interesse, os autores devem declarar: 
``Os autores declaram que não têm nenhum conflito de interesses''. 
Caso contrário, a declaração deve ser: 
``Os autores declaram que têm os seguintes conflitos de interesses: 
lorem ipsum dolor sit amet, consectetur adipiscing elit.''
\end{interests}

\begin{materials}
ESTA DECLARAÇÃO É OBRIGATÓRIA. 
\textit{Considerar o Código de Conduta da SBC. 
Reprodutibilidade de resultados de pesquisa: 
nos casos pertinentes, recomenda-se que artigos indiquem 
a disponibilidade pública de material utilizado nas revisões, 
de modo a facilitar a reprodução dos respectivos resultados por outros pesquisadores.}
%
Se os autores disponibilizaram seus dados e/ou códigos abertamente, a declaração deve ser: 
``Os dados e outros materiais criados e/ou usados nesta revisão da literatura
estão disponíveis em \ldots''. 
Caso contrário, a declaração deve ser: 
``Os dados e outros materiais criados e/ou usados nesta revisão da literatura 
serão disponibilizados mediante solicitação''. 
%
Todo o material suplementar disponibilizado publicamente pelos autores aparecerá 
na página de publicação do artigo em caso de aceitação. 
\end{materials}

\begin{aitools}
ESTA DECLARAÇÃO É OBRIGATÓRIA.  
\textit{Considerar o Código de Conduta da SBC.
Uso de Inteligência Artificial (IA) Generativa: 
a utilização de ferramentas e tecnologias de IA Generativa para geração de conteúdos, 
na escrita e/ou revisão do conteúdo de artigos, deve ser declarada explicitamente no trabalho.} 
A declaração deve ocorrer nesta seção, definida especificamente para este fim.
Deve-se listar as ferramentas (e suas versões) e descrever onde foram empregadas, por exemplo, 
textos, tabelas, gráficos, citações, etc. 
Essas ferramentas \textbf{não} podem ser listadas como autores de um artigo. 
O uso de tais ferramentas não exime os autores da responsabilidade 
sobre todo o seu conteúdo, inclusive no caso de ser identificado plágio.
%
Se os autores não utilizaram tais ferramentas, a declaração deve ser: 
``Os autores não utilizaram ferramentas de IA generativa no desenvolvimento do estudo''.
%
\end{aitools}

\begin{furtherinformation}
Manuscritos devem indicar explicitamente como as questões éticas da pesquisa foram abordadas e gerenciadas, 
incluindo a aprovação da pesquisa por um comitê de ética, quando aplicável.
\end{furtherinformation}

\end{declarations}

\bibliographystyle{apalike-sol}
\bibliography{refs}

\end{document}

